\section{想定される貢献と否定結果}

本稿の現在の貢献は、観測境界、異常分類、入力契約、評価プロトコル、
証拠ゲートを固定することである。今後の実験が仮説を支持した場合、
本研究の経験的貢献は以下になる。

\begin{enumerate}
  \item ストレージ装置内で収集可能な安価な 10 秒 I/O 統計による、
        ransomware anomaly detection の入力契約。
  \item 未知または変種の ransomware を、教師あり family 識別ではなく
        正常時系列からの逸脱として扱う教師なし検出プロトコル。
  \item 正常学習 AE の再構成誤差が、write-heavy ルールや entropy ルールを超えて
        どの条件で有効かを示す評価。
  \item entropy/compression なし、暗号化済みボリューム、benign workload shift に対する
        false positive 分析。
  \item volume あたり 500 KB モデルメモリ制約と、多数 volume 展開時の
        aggregate memory/CPU 制約を同時に見る AE architecture 選択と
        MNN 変換/parity/peak memory 測定手順。
  \item channel-wise error に基づく、ストレージ管理者向け説明可能性の設計。
\end{enumerate}

逆に、AE が単純 baseline を上回らない場合も重要な結果である。その場合は、
未知脅威対応という利点を考慮しても、ストレージ装置内ではルールまたは
別の教師なし軽量検出器が妥当であり、AE は feature discovery または
補助 score に限定すべきである、という結論になる。教師あり軽量分類器が
上回る場合も、それは既知 label が使える運用条件での参照結果として扱い、
未知 family 対応の主張とは分ける。
